\chapter{الكسور أعدادًا}\label{ch:g5:fractions}

\omterm{def:g4:fractions:def}{الكسر} ليس مجرد صورة لفطيرة: بل هو \emph{عدد} حقيقي،
له مكان على مستقيم الأعداد --- وقد يكون المكان نفسه الذي
يقع فيه عدد عشري. يصل هذا الفصل بين العالمين. (أما قواعد
حساب الكسور كاملةً فتأتي في
\cref{ch:g6:fractions} وما بعده.)

\section{الكسور على المستقيم من جديد}

\begin{example}[الوضع والقراءة]\label{ex:g5:fractions:place}
على مستقيم مقطوع أخماسًا، يقع $\frac{7}{5}$ على بُعد $7$ خطوات من $0$:
بعد $1$ (وهو $\frac55$)، عند $1 + \frac25$. فكل \omterm{def:g4:fractions:def}{كسر} يعيش
في مكان ما على المستقيم؛ \omterm{def:g4:fractions:def}{والبسط} الأكبر (عند تساوي \omterm{def:g4:fractions:def}{المقام}) يعني
موضعًا أبعد نحو اليمين.
\end{example}

\begin{proposition}[نقطة واحدة، وأسماء كثيرة]\label{prop:g5:fractions:equal}
قطع كل جزء إلى جزأين (أو ثلاثة\,\dots) لا يحرّك
النقطة:
\[
\frac{1}{2} = \frac{2}{4} = \frac{3}{6} = \frac{5}{10},
\qquad
\frac{2}{3} = \frac{4}{6} = \frac{20}{30}.
\]
فضرب \omterm{def:g4:fractions:def}{البسط} \omterm{def:g4:fractions:def}{والمقام} (أو قسمتهما) على العدد نفسه يعطي اسمًا آخر
\omterm{def:g4:fractions:def}{للكسر} نفسه.
\end{proposition}
\admitted

\begin{omfigure}
\begin{tikzpicture}[scale=0.62]
  \foreach \i in {0,1} \draw[omExo] (\i*4,0) rectangle (\i*4+4,1.0);
  \fill[omDef!40] (0.06,0.06) rectangle (3.94,0.94);
  \draw[thick] (0,0) rectangle (8,1.0);
  \node[right, font=\small] at (8.3,0.5) {$\frac12$};
  \begin{scope}[yshift=-1.7cm]
  \foreach \i in {0,...,3} \draw[omExo] (\i*2,0) rectangle (\i*2+2,1.0);
  \fill[omDef!40] (0.06,0.06) rectangle (1.94,0.94);
  \fill[omDef!40] (2.06,0.06) rectangle (3.94,0.94);
  \draw[thick] (0,0) rectangle (8,1.0);
  \node[right, font=\small] at (8.3,0.5) {$\frac24$};
  \end{scope}
  \begin{scope}[yshift=-3.4cm]
  \foreach \i in {0,...,9} \draw[omExo] (\i*0.8,0) rectangle
    (\i*0.8+0.8,1.0);
  \foreach \i in {0,...,4} \fill[omDef!40] (\i*0.8+0.05,0.06)
    rectangle (\i*0.8+0.75,0.94);
  \draw[thick] (0,0) rectangle (8,1.0);
  \node[right, font=\small] at (8.3,0.5) {$\frac{5}{10}$};
  \end{scope}
\end{tikzpicture}

{\small ثلاثة أشرطة، وثلاثة أسماء، ومقدار واحد: $\frac12 = \frac24 =
\frac{5}{10}$.}
\end{omfigure}

\section{الكسور والأعداد العشرية}

\begin{proposition}[كسور لها أسماء عشرية]\label{prop:g5:fractions:decimal}
كسور الأعشار والأجزاء من مئة والأجزاء من ألف \emph{هي} أعداد
عشرية:
\[
\frac{3}{10} = 0.3, \qquad \frac{27}{100} = 0.27 .
\]
وبعض الكسور الأخرى له أسماء عشرية أيضًا، تُوجد بالانتقال إلى
الأعشار أو الأجزاء من مئة:
\[
\frac{1}{2} = \frac{5}{10} = 0.5, \qquad
\frac{1}{4} = \frac{25}{100} = 0.25, \qquad
\frac{3}{4} = 0.75, \qquad
\frac{1}{5} = \frac{2}{10} = 0.2 .
\]
لكن ليس كلها: \omterm{def:g4:fractions:def}{فالكسر} $\frac13 = 0.333\dots$ لا ينتهي أبدًا ---
\omterm{def:g4:fractions:def}{والكسر} هو اسمه الدقيق الوحيد.
\end{proposition}
\admitted

\begin{omfigure}
\begin{tikzpicture}[scale=1.55]
  \draw[-Stealth] (-0.2,0) -- (7.6,0);
  \foreach \i in {0,...,28} \draw ({\i*0.25},-0.05) -- ({\i*0.25},0.05);
  \foreach \i in {0,4,...,28} \draw ({\i*0.25},-0.09) -- ({\i*0.25},0.09);
  \foreach \x in {0,...,7} \node[below, font=\small] at (\x,-0.1) {$\x$};
  \fill[omDef] (0.5,0) circle (1.6pt)
    node[above, font=\small] {$\frac12 = 0.5$};
  \fill[omThm] (2.75,0) circle (1.6pt)
    node[above, font=\small] {$\frac{11}{4} = 2.75$};
  \fill[omMeth] (6.25,0) circle (1.6pt)
    node[above, font=\small] {$\frac{25}{4} = 6.25$};
\end{tikzpicture}

{\small مستقيم واحد، ولغتان: فكل نقطة يمكن أن تحمل اسمًا كسريًّا
واسمًا عشريًّا.}
\end{omfigure}

\begin{example}[اختيار الاسم الأيسر]\label{ex:g5:fractions:choose}
لدفع ثلاثة أرباع $12$: الاسم الكسري أيسر،
إذ $\frac34$ من $12$ هو $9$. ولمقارنة $\frac34$ و $0.8$: الاسم
العشري أيسر، $0.75 < 0.8$. فإتقان اللغتين يفيد.
\end{example}

\section{مقارنة الكسور}

\begin{method}[ثلاث أدوات للمقارنة]\label{met:g5:fractions:compare}
\begin{enumerate}
  \item \emph{\omterm{def:g4:fractions:def}{المقام} نفسه}: الأجزاء الأكثر تفوز،
        $\frac57 > \frac37$؛
  \item و\emph{\omterm{def:g4:fractions:def}{البسط} نفسه}: الأجزاء الأكبر تفوز، \omterm{def:g4:fractions:def}{فالمقام}
        \emph{الأصغر} يفوز: $\frac35 > \frac38$ (فالأخماس أكبر
        من الأثمان)؛
  \item و\emph{المعالم}: قارن كلًّا منهما بالعدد $\frac12$ أو بالعدد $1$:
        $\frac38 < \frac12 < \frac59$، إذن $\frac38 < \frac59$.
\end{enumerate}
\end{method}

\begin{example}\label{ex:g5:fractions:compare}
أيّهما أكل بيتزا أكثر: علي وقد أكل $\frac58$ أم بثينة وقد أكلت
$\frac54$؟ \omterm{def:g4:fractions:def}{كسر} بثينة أكبر من $1$ (فخمسة أرباع أكثر من بيتزا
كاملة واحدة)، \omterm{def:g4:fractions:def}{وكسر} علي أصغر من $1$: فبثينة أكلت أكثر --- بل
أكثر من بيتزا كاملة.
\end{example}

\section{تمارين}

\begin{exercise}[$\star$]\label{exo:g5:fractions:1}
اُرسم مستقيم أعداد مقطوعًا أخماسًا من $0$ إلى $2$، وضع عليه:
$\frac25$؛ \ $\frac55$؛ \ $\frac85$؛ \ $\frac{10}{5}$.
\end{exercise}

\begin{exercise}[$\star$]\label{exo:g5:fractions:2}
أكمل الأسماء المتساوية:
$\frac12 = \frac{?}{8}$; \quad
$\frac34 = \frac{?}{8}$; \quad
$\frac{2}{5} = \frac{4}{?}$; \quad
$\frac{6}{9} = \frac{2}{?}$.
\end{exercise}

\begin{exercise}[$\star$]\label{exo:g5:fractions:3}
أي الكسور تسمّي أعدادًا صحيحة؟ أعطِ العدد الصحيح عند ذلك:
$\frac82$؛ \quad $\frac{9}{4}$؛ \quad $\frac{15}{3}$؛ \quad
$\frac{20}{5}$.
\end{exercise}

\begin{exercise}[$\star$]\label{exo:g5:fractions:4}
اُكتب على صورة أعداد عشرية: $\frac{7}{10}$؛ \quad $\frac{41}{100}$؛ \quad
$\frac12$؛ \quad $\frac34$؛ \quad $\frac15$.
\end{exercise}

\begin{exercise}[$\star$]\label{exo:g5:fractions:5}
اُكتب على صورة كسور (أعشار أو أجزاء من مئة): $0.9$؛ \quad $0.13$؛
\quad $2.5$؛ \quad $0.75$ (ثم أعطِ اسمه بالأرباع).
\end{exercise}

\begin{exercise}[$\star$]\label{exo:g5:fractions:6}
قارن، مع تسمية الأداة المستعملة
(\cref{met:g5:fractions:compare}):
\[
\frac47 \;?\; \frac67, \qquad
\frac35 \;?\; \frac37, \qquad
\frac25 \;?\; \frac{7}{12}, \qquad
\frac98 \;?\; \frac{11}{12} .
\]
\end{exercise}

\begin{exercise}[$\star$]\label{exo:g5:fractions:7}
رتّب بالأسماء العشرية: $\frac12$؛ \ $0.4$؛ \ $\frac34$؛ \
$0.6$؛ \ $\frac15$.
\end{exercise}

\begin{exercise}[$\star$]\label{exo:g5:fractions:8}
ربع تلاميذ صفّ من $28$ يعزفون آلة، ونصفهم يمارسون
رياضة. فكم تلميذًا في كل مجموعة؟ وأي المجموعتين أكبر؟
\end{exercise}

\begin{exercise}[$\star$]\label{exo:g5:fractions:9}
اِربط كل \omterm{def:g4:fractions:def}{كسر} بمدة: $\frac14$ سا، و $\frac12$ سا،
و $\frac34$ سا، و $\frac{1}{60}$ سا --- والمدد هي $30$ د، و $45$ د، و $1$ د،
و $15$ د.
\end{exercise}

\begin{exercise}[$\star\star$]\label{exo:g5:fractions:10}
شرب عمر $\frac35$ من قنينته \omterm{def:g2:measure:units}{سعة} $50$ سل؛ وشربت آية $\frac12$ من
قنينتها \omterm{def:g2:measure:units}{سعة} $60$ سل. فأيّهما شرب أكثر؟ (احسب المقدارين بالسنتيلتر ---
فمقارنة الكسور المجرّدة لا تكفي هنا، وقل لماذا.)
\end{exercise}

\begin{exercise}[$\star\star$]\label{exo:g5:fractions:11}
أعطِ كسرًا يقع تمامًا بين $\frac12$ و $1$؛ ثم كسرًا يقع
تمامًا بين $\frac12$ و $\frac34$. (وقد تساعد الأسماء
العشرية.)
\end{exercise}
